\section{Introduction}
\label{sec:intro}

In the tasks studied here, background texture changes neither the requested goal nor the object's
dynamics or reward. A useful abstraction can ignore that texture while preserving where the object
is and what an action can do to it. Joint-embedding predictive architectures (JEPAs) learn this
abstraction by predicting future representations rather than reconstructing every future pixel
\citep{lecun2022path,assran2023ijepa,bardes2024vjepa,maes2026lewm}. The promise is attractive for
physical intelligence: learn a compact description of the world that can be reused across goals.
The difficulty is that predictive learning must decide both \emph{what to represent} and
\emph{what kind of prediction a decision needs}.

These are separate requirements. First, an easy-to-predict feature need not be useful for control.
A background fixed within an episode is perfectly predictable even if the controlled object is
ignored \citep{sobal2022slow}. Second, retaining the right state does not make its mean future a
sufficient description of uncertain outcomes. An action that lands equally often on either side
of a target has its mean exactly on the target. A planner that scores that mean can prefer it to
an action that reliably lands close by. The mean prediction can be correct and the decision wrong.

We study both requirements in a minimal visual world: one red cube, four movement actions, a
background, and optional stochastic wind or a moving wall. The simplicity is deliberate. Cube
position, the requested goal, and the executed path are directly visible, so we can distinguish
three cases without relying on a large game's aggregate score:
\begin{enumerate}
\item \textbf{Readable state and goal-directed movement.} Deterministic movement with changing
visual noise is the initial control condition. Requested visual goals and actual trajectories test
planning beyond merely staying alive.
\item \textbf{Failing to expose the controlled state.} Varying background persistence and wind tests
whether the same predictive-learning recipe retains an accessible description of the cube.
\item \textbf{Retaining state while predicting an uncertain future.} Stochastic movement and a
late-revealed opening test the decisions made from a learned imagination model. Exact finite-branch
calculations separately isolate the effect of replacing a distribution by its mean.
\end{enumerate}

Two elementary calculations make the third case precise. The first separates the cost of a mean
from the expected cost of an outcome. The second shows that recursively applying an exact one-step
conditional mean need not predict a later mean. Both hold under ideal prediction, without an
optimization failure. A Gaussian nuisance construction supplies the corresponding limit on the
representation-learning objective. These mechanisms motivate the experiments; the experiments
determine which failure is actually observed in each learned model.

The distinctive empirical result is a reversal in the effect of nuisance persistence. Under
predictable motion, changing gray noise yields more accurate position read-out than persistent noise;
with fresh wind innovations, changing noise fails while persistent noise preserves subpixel read-out.
The architecture, latent dimension, training data size, updates, and three-model replication are
matched across this comparison. We use actual goal-directed paths and a shared-policy JEPA/flow
comparison to determine what this representation boundary means for behavior. The point is not to
rediscover that prediction and control differ, but to identify when the same predictive-learning
recipe changes which physical state it exposes.
